\documentclass[]{report}

\usepackage{import}

\import{../}{settings.tex}

\title{Cours d'intégration pour l'agrégation}

\author{Cours de David GERARD-VARET, transcrit par SADR TAHOURI Luca}
\date{}

\begin{document}

    \maketitle
    \tableofcontents

\chapter{Généralités sur les groupes}

\section{Définitions,premières propriétés,exemples}

\begin{df}{Groupe}{}
    Soit $G$ un ensemble et $*$ une lois de composition interne.\\
    $G$ est un groupe si:
    \begin{enumerate}[label=\roman*]
        \item 
    \end{enumerate}
\end{df}

\begin{re}{}{}
    Le neutre $e_G$ est unique ($e*e'=e=e'$).\\
    On peut toujours "simplifier" par un élément dans un groupe
\end{re}

\begin{df}{Groupe commutatif}{}
    Si la loi de $(G,*)$ est commutative, on dit que $G$ est abélien:
    \[\forall x,y \in G,x*y=y*x\]
    On réserve la notation "$+$" aux groupes abéliens.
\end{df}

\begin{df}{Ordre d'un groupe}{}
    L'ordre d'un groupe est son cardinal
\end{df}

\begin{df}{}{}
    Soit $X$ un ensemble. Une permutation de $X$ est une bijection de $X$ dans $X$.
    \[\Sr_X=\{\text{permutations de $X$}\}\text{ est un groupe pour $\circ$}\]
\end{df}

\begin{df}{}{}
    Soit $G$ un groupe, une partie $H$ de $G$ est un sous-groupe de $G$ si et seulement si muni de la restriction de la loi sur $G$, $H$ a une structure de groupe
\end{df}

\begin{prop}{}{}
    $G$ groupe, $H\subset G$. $H$ est un sous-groupe si:
    \begin{itemize}
        \item $H\neq \emptyset$
        \item $\forall x,y \in H,x y^{-1}\in H$
    \end{itemize}
\end{prop}

\begin{prop}{}{}
    L'intersection de 2 sous-groupes est un sous-groupe.
\end{prop}

\begin{df}{Sous-groupe engendré}{}
    Soit $G$ un groupe, $S\subset G$ une partie quelconque. Le sous groupe engendré par $S$, noté $<S>$ est l'intersection de tous les sous groupes de $G$ contenant $S$.\\
    C'est le plus petit sous-groupe de $G$ contenant $S$
\end{df}

\begin{prop}{}{}
    \[<S>=\{ s_1^{\alpha_1}\ldots s_n^{\alpha_n},n\in \N,\alpha_i \in \Z,s_i\in S\}\]
\end{prop}

\begin{ex}{Groupe diédral $D_n$}{}
    $n\in \N,n\ge 2$. $D_n$ est le groupe des isométries du plan qui laissent invariant un polygone régulier à $n$ cotés.
\end{ex}

\begin{df}{Ordre d'un élément d'un groupe}{}
    $G$ un groupe, $g\in G$, l'ordre de $g$ est le cardinal de $<g>$
\end{df}

\begin{prop}{}{}
    Soit $G$ un groupe, $g\in G$. Alors, $\forall k\in \N, g^k=e\iff $ ordre de $g$ divise $k$.\\
    Si $g$ est l'ordre de $g$, alors $g^k$ est d'ordre.
    \[\frac{d}{\pgcd(d,k)}\]
\end{prop}

\section{Classes à gauche, à droite, théorème de Lagrange}

\begin{prop}{}{}
    $\forall g\in G$
    \[\fundf{H}{gH}{x}{gx}\text{ est une bijection}\]
\end{prop}

\begin{cor}{}{}
    Comme $G=\sqcup_{\overline{g}\in G/H} gH$, 
    \[\left|G\right|=\sum_{\overline{g}\in G/H} \left|H\right|= \left|G/H\right||H|\]
\end{cor}

\begin{tm}{Théorème de Lagrange}{}
    Si $G$ est fini, $H$ est un sous-groupe de $G$, alors $|H|$ divise $G$
\end{tm}

\begin{cor}{}{}
    Si $|G|=n$, alors $\forall g\in G, g^n=e$
\end{cor}

\begin{df}{}{}
    $[G:H]$ est l'indice de $H$ dans $G$, et est définit par $|G/H|=|H\wt G|$.\\
    Si $G$ est fini, $[G:H]=\frac{|G|}{|H|}$
\end{df}

\begin{df}{}{}
    Le produit direct de deux groupes $G,G'$ est le produit cartésien $G\times G'$ muni des produits composante par composante. $\forall (x,y,x',y')\in G^2\times G'^2$,
    \[(x,x')*(y,y')=(x*y,x'*y')\]
    On peut le faire avec un nombre fini quelconque de groupes
\end{df}

\section{Morphismes, isomorphismes}

\begin{df}{Morphisme de groupes}{}
    A VOIR
\end{df}

\begin{re}{}{}
    $\funto{\varphi}{G}{G'}$ morphisme, alors $\varphi(e_G)=e_{G'}$
\end{re}

\begin{df}{Image directe, image réciproque d'un sous-groupe, Isomorphisme, Automorphisme}{}
    A VOIR
\end{df}

\begin{prop}{}{}
    Aut$(G)$ l'ensemble des automorphismes, d'un groupe $G$ est un groupe pour $\circ$.
\end{prop}

\begin{df}{Noyau d'un morphisme de groupes}{}
    A VOIR
\end{df}

\begin{prop}{}{}
    $\varphi$ est injective $\iff \ker \varphi=\{e\}$.
\end{prop}

\section{Sous-groupes distingués, groupe quotient}

$G$ un groupe et $H$ un sous-groupe de $G$. On aimerait mettre sur $G/H$ (ou $H\wt G$) une structure de groupe par $g_1 H*g_2 H=g_1 g_2 H$ (*).\\
C'est demander que la surjection canonique soit un morphisme de groupes et ce n'est pas toujours possible !\\
Condition nécessaire, que (*) ne dépende pas du choix des représentants dans les classes:
\[\forall g_1,g_2\in G,\forall h_1,h_2\in H, g_1 g_2=g_1 h_1 g_2 h_2 H\]
c'est à dire
\[g_2 H=h_1g_2 H\]
En particulier, $h_1 g_2\in g_2 H$. C'est vrai $\forall h_1\in H$, donc $Hg_2\subset g_2 H$. Comme c'est vrai pour tout $g_2\in G$, et que les classes à gauche, à droite forment une partition de $G$, on a
\[\forall g\in G ,gH=Hg\]
On peut vérifier que cette condition est suffisante.

\begin{df}{Groupe distingué}{}
    Soit $G$ un groupe, $H$ un sous-groupe de $G$, on dit que $H$ est distingué (ou normal) dans $G$ noté $H\triangleleft G$ si $\forall g\in G, gHg^{-1}=H$.
\end{df}

\begin{prop}{}{}
    $G$ un groupe, $H$ un sous-groupe de $G$.
    \begin{enumerate}
        \item Si $H\triangleleft G$, alors $G/H$ a une (unique) structure de groupe tel que la surjection canonique est un morphisme de groupes.
        \item Réciproquement, si $G/H$ a une structure de groupes tq la surjection canonique est un morphisme, alors $H\triangleleft G$
    \end{enumerate}
\end{prop}

\begin{re}{}{}
    Dans un groupe abélien, tous les sous-groupes sont distingués.
\end{re}

\begin{ex}{}{}
    $\Z/n\Z$ est un groupe
\end{ex}

\begin{prop}{}{}
    $\funto{\varphi}{G}{G'}$ un morphisme de groupe, alors $\ker \varphi \triangleleft G$
\end{prop}

\begin{proof}
    Soit $x\in \ker (\varphi)$, soit $g\in G$, alors
    \begin{align*}
        \varphi (g x g^{-1})&=\varphi(g)\varphi(x)\varphi(g)^{-1}\\
        &=\varphi (g) e_{G'} \varphi(g)^{-1}\\
        &=e_{G'}
    \end{align*}
\end{proof}

\begin{prop}{}{}
    Si $H\triangleleft G$, alors $H$ est le noyau d'un morphisme de $G$ vers un groupe.
    \[H=\ker (\pi)\]
    ou $\pi$ est la surjection canonique
\end{prop}
\end{document}